\section{Công việc liên quan}
\label{sec:relatedWork}

Trước tiên, chúng tôi thảo luận về công việc liên quan về phát hiện và giải thích bất thường theo ngữ cảnh, sau đó tiến hành hai loại phát hiện bất thường truyền thống có liên quan chặt chẽ nhất: phát hiện bất thường dựa trên phụ thuộc và dựa trên không gian con.

\subsection{Phát hiện và giải thích sự bất thường theo ngữ cảnh}

Phát hiện bất thường theo ngữ cảnh đã nhận được sự chú ý đặc biệt trong dữ liệu không gian \citep{cai2013spatial}, dữ liệu thời gian \citep{salvador2004learning} và dữ liệu không gian-thời gian \citep{smets2009discovering}, trong đó các đặc điểm không gian và/hoặc thời gian được sử dụng để xác định bối cảnh. Các phương pháp này không thể áp dụng trực tiếp cho các miền khác, trong đó bối cảnh được xác định bởi các loại đối tượng địa lý khác; Phát hiện bất thường theo ngữ cảnh 'chung' đã nhận được sự quan tâm hạn chế trong cộng đồng.

CAD \citep{song2007conditional} là một công trình quan trọng giới thiệu tính năng phát hiện bất thường theo ngữ cảnh chung. Nó giả định một phân vùng các tính năng do người dùng chỉ định thành các đặc trưng ngữ cảnh (được gọi là 'môi trường') và hành vi (được gọi là 'chỉ báo'), đồng thời sử dụng Mô hình hỗn hợp Gaussian để phù hợp với sự phân bố của các không gian đặc trưng ngữ cảnh và hành vi. Sự phụ thuộc giữa các đặc điểm ngữ cảnh và đặc trưng hành vi sau đó được học bằng 'các hàm ánh xạ' và một đối tượng được coi là bất thường nếu nó vi phạm các hàm đã học. Để điều này hoạt động, CAD giả định rằng cả đặc điểm ngữ cảnh và hành vi đều bao gồm một số lượng không xác định của nhiều thành phần Gaussian, đây có thể là một giả định chắc chắn trong thực tế. Hơn nữa, CAD chỉ có thể xử lý các tính năng số và rất tốn kém về mặt tính toán. Phương pháp được đề xuất của chúng tôi không đưa ra giả định nào về việc phân phối các tính năng, có thể xử lý các đặc trưng ngữ cảnh hỗn hợp và các đặc trưng hành vi số và chúng tôi sẽ chứng minh bằng thực nghiệm rằng nó hiệu quả hơn về mặt tính toán so với CAD trong khi đạt được độ chính xác phát hiện cao hơn.

ROCOD \citep{liang2016robust} cũng có liên quan chặt chẽ và sử dụng các mô hình hành vi dự kiến ​​cục bộ và toàn cầu để mô tả sự phụ thuộc giữa các đặc trưng ngữ cảnh và hành vi. Cụ thể, các mô hình hồi quy tiêu chuẩn như CART được sử dụng để tìm hiểu các mô hình toàn cầu, với các đặc trưng ngữ cảnh là các biến dự đoán và các đặc trưng hành vi là các biến phản hồi. Các mẫu cục bộ được tính toán dựa trên phương tiện giá trị đặc trưng hành vi của các láng giềng của đối tượng. Giá trị thực tế của một đối tượng được so sánh với mô hình cục bộ và tổng thể, và mức trung bình có trọng số của những khác biệt này tạo thành điểm bất thường. Vì giá trị trung bình có điều kiện chỉ mô tả một khía cạnh của phân bố có điều kiện và không nhất thiết là điểm có xác suất xảy ra cao nhất (tức là chế độ). Để giải quyết vấn đề này, phương pháp của chúng tôi sử dụng phân tích hồi quy phân vị để ước tính các phân vị có điều kiện, cung cấp mô tả đầy đủ hơn về phân bố có điều kiện. Kết quả là, phương pháp của chúng tôi về mặt thực nghiệm vượt trội hơn ROCOD về độ chính xác.Với việc sử dụng ngày càng nhiều các thuật toán phát hiện sự bất thường trong các lĩnh vực quan trọng về an toàn, chẳng hạn như chăm sóc sức khỏe và sản xuất, nghĩa vụ đạo đức và quy định trong việc đưa ra lời giải thích cho các quyết định mang tính rủi ro cao do các thuật toán này đưa ra đã trở nên rõ ràng hơn \citep{li2022survey}. Tuy nhiên, tính năng phát hiện bất thường theo ngữ cảnh hiện tại không đưa ra những giải thích như vậy và sẽ cần có những sửa đổi không hề nhỏ để thực hiện điều đó. Cụ thể, CAD \citep{song2007conditional} tính toán điểm bất thường của một phiên bản dữ liệu bằng cách đo độ lệch của nó so với các hàm ánh xạ được học từ phần lớn các phiên bản dữ liệu. Cách tiếp cận này đặt ra thách thức trong việc tạo ra các giải thích nội tại, vì nó không cho phép quy kết điểm bất thường cho các đặc điểm riêng lẻ. Hơn nữa, ROCOD \citep{liang2016robust}, một phương pháp được biết đến khác để phát hiện điểm bất thường theo ngữ cảnh, tính toán mức trung bình có trọng số của sự khác biệt của một phiên bản với các mẫu cục bộ và toàn cầu đã học dưới dạng điểm bất thường của nó. Mẫu cục bộ được lấy bằng cách sử dụng các mẫu lân cận của nó, trong khi mẫu chung được học bằng cách sử dụng tất cả các trường hợp, khiến việc liên kết điểm bất thường với các đặc điểm riêng lẻ trở nên khó khăn.

Bất chấp sự tồn tại của nhiều phương pháp để giải thích các dị thường, chẳng hạn như những phương pháp được nêu trong các bài báo (khảo sát) gần đây \citep{panjei2022survey,li2022survey,xu2021beyond}, vẫn có rất ít nghiên cứu về cách giải thích các bất thường theo ngữ cảnh. Đặc biệt, COIN \citep{liu2018contextual} giải thích các ngoại lệ bằng cách báo cáo điểm ngoại lệ của chúng, các đặc điểm góp phần tạo ra sự bất thường của nó và mô tả theo ngữ cảnh của các vùng lân cận. Tuy nhiên, COIN xử lý tất cả các tính năng như nhau, trong khi phương pháp của chúng tôi chia các tính năng thành các đặc trưng ngữ cảnh và hành vi. Hơn nữa, chúng tôi phát triển một hình ảnh trực quan giúp giải thích các điểm bất thường theo ngữ cảnh bên cạnh việc báo cáo điểm bất thường tổng thể, tầm quan trọng của đối tượng địa lý, các lân cận theo ngữ cảnh và điểm bất thường riêng lẻ cho từng đặc trưng hành vi.



Các phương pháp sau đây ít liên quan hơn vì chúng xem xét các vấn đề hơi khác nhau. \cite{valko2011conditional} xây dựng một phương pháp dựa trên biểu đồ phi tham số để phát hiện sự bất thường có điều kiện, phương pháp này chỉ giải quyết vấn đề về một đặc trưng hành vi phân loại duy nhất. \cite{hayes2014contextual} trước tiên xác định các điểm bất thường trên các đặc trưng hành vi, sau đó tinh chỉnh các điểm bất thường được phát hiện bằng cách phân cụm tất cả các đối tượng trên các đặc trưng ngữ cảnh. \cite{tang2015mining} phát hiện các điểm bất thường của nhóm từ dữ liệu phân loại đa chiều.  \cite{hong2015multivariate} trình bày một khung phát hiện bất thường theo ngữ cảnh dành riêng cho các đặc trưng hành vi phân loại, khung này cũng xem xét sự phụ thuộc giữa các đặc trưng hành vi. Hơn nữa, \cite{zheng2017contextual} áp dụng phương pháp học số liệu mạnh mẽ trên các đặc trưng ngữ cảnh để tìm các lân cận theo ngữ cảnh có ý nghĩa hơn và sau đó tận dụng hồi quy hạt nhân $k$-NN để dự đoán các giá trị đặc trưng hành vi. \cite{meghanath2018conout} phát triển ConOut để tự động tìm và kết hợp nhiều bối cảnh nhằm xác định và giải thích các ngoại lệ.


\subsection{Phát hiện bất thường truyền thống}Mặc dù phát hiện bất thường truyền thống coi một vấn đề khác với phát hiện bất thường theo ngữ cảnh, các kỹ thuật tận dụng phát hiện bất thường dựa trên phụ thuộc và dựa trên không gian con có liên quan đến phương pháp của chúng tôi.

\subsubsection{Phát hiện bất thường dựa trên phụ thuộc}
Tính năng phát hiện bất thường dựa trên phần phụ thuộc nhằm mục đích xác định các điểm bất thường bằng cách khám phá các phần phụ thuộc giữa các tính năng. \cite{teng1999correcting} khám phá sự phụ thuộc giữa các tính năng không phải mục tiêu và mục tiêu để xác định và sửa các điểm dữ liệu nhiễu có thể xảy ra. Để phát hiện các hành vi xâm nhập mạng, \cite{huang2003cross} trình bày Phân tích tính năng chéo để nắm bắt các mẫu phụ thuộc giữa các tính năng trong lưu lượng mạng thông thường. Để phát hiện sự bùng phát dịch bệnh, \cite{wong2003bayesian} đề xuất khám phá sự phụ thuộc giữa các tính năng bằng mạng Bayesian. \cite{noto2010anomaly} đề xuất phát hiện các điểm bất thường bằng cách sử dụng một tập hợp các mô hình, trong đó mỗi mô hình khám phá sự phụ thuộc giữa tính năng phản hồi và các tính năng khác. \cite{babbar2012mining} sử dụng mạng Bayesian Gaussian tuyến tính để phát hiện những điểm bất thường vi phạm mối quan hệ nhân quả.

LoPAD \citep{lu2020lopad} trước tiên sử dụng mạng Bayesian để tìm Chăn Markov của từng tính năng. Sau đó, một mô hình dự đoán (ví dụ: GIỎ HÀNG) được học cho từng tính năng riêng lẻ dưới dạng biến phản hồi, với Chăn Markov của nó làm biến dự đoán. Với một đối tượng, LoPAD tính toán khoảng cách Euclide giữa giá trị thực tế và giá trị dự đoán (tức là giá trị trung bình có điều kiện) làm điểm bất thường cho từng đối tượng. Điểm bất thường thu được được chuẩn hóa và tính tổng để có được điểm bất thường cuối cùng cho một đối tượng. Phương pháp này không phân biệt các đặc trưng ngữ cảnh và hành vi và --- giống như ROCOD --- sử dụng các phương tiện có điều kiện để biểu diễn phân bố có điều kiện. Do đó, LoPAD không thể (chính xác) phát hiện các điểm bất thường theo ngữ cảnh.

\subsubsection{Phát hiện bất thường dựa trên không gian con}
Phát hiện dị thường dựa trên không gian con tìm cách tìm ra các dị thường trong một phần của không gian đặc trưng. Cụ thể, \cite{kriegel2009outlier} đề xuất SOD để xác định các ngoại lệ trong các không gian con khác nhau. Cụ thể, họ xây dựng các không gian con trục song song được kéo dài bởi các lân cận của một đối tượng nhất định. Trên cơ sở này, họ điều tra xem liệu đối tượng này có sai lệch đáng kể so với các đối tượng lân cận của nó trên bất kỳ không gian con nào trong số này hay không. Hơn nữa, \cite{kriegel2012outlier} mở rộng công việc này để xác định xem một đối tượng có bất thường trên các không gian con được định hướng tùy ý bao quanh các không gian lân cận của nó hay không. \cite{nguyen2013cmi} đề xuất tìm các không gian con có mối tương quan chặt chẽ với nhau và sau đó xác định các điểm bất thường trên các không gian con này. Cuối cùng, \cite{cabero2021archetype} sử dụng phân tích nguyên mẫu để chiếu không gian đặc trưng vào các không gian con khác nhau với các mối tương quan tuyến tính dựa trên các lân cận gần nhất. Trên cơ sở này, họ khám phá các giá trị ngoại lệ bằng cách tập hợp các kết quả thu được trên các không gian con có liên quan. Nhìn chung, các phương pháp này nhằm mục đích xác định các điểm bất thường trong một tập hợp con các đặc điểm, nhưng xử lý tất cả các đặc điểm một cách bình đẳng và do đó không phù hợp để xác định các điểm bất thường theo ngữ cảnh.